% !TeX root = ../机器学习笔记.tex
% 由主文档通过 \input 载入；请编译主文档。
\setcounter{chapter}{11} % 保留原书章号。
\chapter{集成学习}
\label{chap:ml-ensembles}

集成学习（ensemble learning）组合多个基学习器。平均试图削弱个体波动，提升则逐步纠正已有模型的不足；两种思想都依赖组合成员之间的关系，而非只依赖模型数量。

\section{模型组合、差异性与误差相关}

对回归预测 $f_1,\ldots,f_M$，平均预测为 $\bar f=M^{-1}\sum_m f_m$。若在固定输入处每个预测具有相同方差 $\sigma^2$，两两协方差为 $\rho\sigma^2$，则
\[
 \operatorname{Var}(\bar f)=\sigma^2\left(\rho+\frac{1-\rho}{M}\right).
\]
这由方差展开的 $M$ 个对角项与 $M(M-1)$ 个非对角项直接得到。独立误差平均后方差降为 $\sigma^2/M$，完全相同的预测则没有这种收益。共同偏差不会因增加同类模型自动消失。

分类可作多数投票或平均预测概率；投票只保留类别，概率平均还保留置信信息。加权组合必须说明权重怎样选择，若利用测试表现选择权重，测试集便参与了训练流程。

\section{自助聚合与袋外评价}

自助聚合（bootstrap aggregating，Bagging）从 $n$ 个训练样本中有放回抽取 $n$ 次构成一个自助样本，训练一个模型，重复 $M$ 次后平均或投票。某原始样本没有进入一次抽样的概率为 $(1-1/n)^n$，大样本时接近 $e^{-1}$。这些未进入该模型训练的样本称袋外（out-of-bag）样本。

对每个样本，只使用没有训练过它的基学习器进行袋外预测，再汇总指标。这是一种有用的内部评价方式，但不能免除预处理和调参引起的信息复用。若各基学习器极其稳定，重采样产生的差异很小，Bagging 的改善也可能有限。

\section{随机森林}

随机森林（random forest）在样本重采样之外，每个树节点只从随机抽取的一部分特征中选择切分。强预测特征不再支配所有树，往往能降低树间相关性。预测时回归取平均，分类按规定采用投票或概率平均。

树数增加主要降低有限随机集成的波动，树深和最小叶样本数控制单树复杂度，候选特征数则影响单树能力与相关性。较多树不等于无需验证，训练数据偏差、泄漏和不适合的目标仍会影响结果。概率输出也可能需要独立校准。

\section{AdaBoost 与指数损失}

二分类标签为 $\pm1$，弱分类器 $h_t(x)\in\{-1,+1\}$。AdaBoost（adaptive boosting）初始化 $D_1(i)=1/n$，按当前权重训练 $h_t$，计算错误率 $\epsilon_t=\sum_iD_t(i)\mathbf1\{h_t(x_i)\ne y_i\}$。若 $0<\epsilon_t<1/2$，置
\begin{equation}
 \alpha_t=\frac12\log\frac{1-\epsilon_t}{\epsilon_t},\qquad
 D_{t+1}(i)=\frac{D_t(i)e^{-\alpha_ty_ih_t(x_i)}}{Z_t}.
 \label{eq:ml-adaboost-update}
\end{equation}
其中 $Z_t$ 使权重和为一。最终按 $F_T(x)=\sum_t\alpha_th_t(x)$ 的符号分类。

该系数最小化 $(1-\epsilon_t)e^{-\alpha}+\epsilon_te^\alpha$，因为其导数在上述 $\alpha_t$ 为零。此时 $Z_t=2\sqrt{\epsilon_t(1-\epsilon_t)}$，而指数训练损失满足 $n^{-1}\sum_i e^{-y_iF_T(x_i)}=\prod_tZ_t$；误分类样本的指数损失至少为一，故训练错误率受该乘积控制。这不等于测试错误保证。若 $\epsilon_t=0$ 可终止或限制系数，等于 $1/2$ 时无改进，大于 $1/2$ 时需换学习器或在允许时翻转预测；噪声标签可能获得越来越大的权重。

\section{梯度提升与二阶近似}

梯度提升（gradient boosting）构造加法模型 $F_t=F_{t-1}+\nu\rho_th_t$，其中 $0<\nu\leq1$ 是收缩率。先从最佳常数模型开始，计算负梯度
$r_{it}=-\partial\ell(a,y_i)/\partial a|_{a=F_{t-1}(x_i)}$，
再用基学习器拟合 $r_{it}$，通过线搜索或叶内优化选 $\rho_t$。平方损失 $\ell(a,y)=(a-y)^2/2$ 时，负梯度就是残差 $y_i-F_{t-1}(x_i)$。

对于二次可微损失，记当前一、二阶导数为 $g_i,h_i$。固定树结构、叶 $R$ 输出常数 $c$，加二次惩罚 $\lambda c^2/2$ 的局部近似为
$c\sum_{i\in R}g_i+\frac12c^2(\sum_{i\in R}h_i+\lambda)$。
分母正时，最佳近似叶值为 $-\sum g_i/(\sum h_i+\lambda)$。这是二阶近似的解，不自动等于原损失的全局最优更新；树深、学习率和轮数需要联合验证。

\section{例题、练习与学习检查}

\begin{example}[提升如何重分配注意力]
四个等权样本中一个被弱分类器误判。
\end{example}
\begin{solve}
$\epsilon=1/4$，$\alpha=\frac12\log3$。误判点权重乘 $\sqrt3$，其余点乘 $1/\sqrt3$，归一后误判点权重为 $1/2$，三个正确点各为 $1/6$。算法使下一轮更重视当前难点，但难点若是错标也会被放大。
\end{solve}

习题：证明本章平均方差公式；推导 AdaBoost 的 $Z_t$；对响应 $0,0,2,2$ 从均值模型开始，用一个树桩作一次平方损失提升。堆叠（stacking）若训练第二层模型，必须用第一层的折外预测构造训练输入，不能用过度拟合的训练内预测代替。





